\documentclass[aspectratio=169,10pt,english]{beamer}

\input{../../../_preambulo_beamer.tex}
\input{../../../_comandos_beamer.tex}

\selectlanguage{english}

\title{MA1001B - Statistical Modeling for Decision Making}
\subtitle{Lesson 25: Sampling Distribution of a Proportion}
\author{}
\institute{Tecnol\'ogico de Monterrey}
\date{}

\begin{document}

\begin{frame}[plain]
  \vspace{1.2cm}
  {\Large\bfseries MA1001B - Statistical Modeling for Decision Making\par}
  \vspace{0.7cm}
  {\large Lesson 25: Sampling Distribution of a Proportion\par}
  \vspace{0.7cm}
  {\color{TecRojo}\rule{\textwidth}{1.2pt}}
  \vspace{0.8cm}

  MA1001B Faculty\\[0.3cm]
  Term\\[0.3cm]
  Tecnol\'ogico de Monterrey
\end{frame}

\begin{frame}{The Sample-Proportion Question}
  \begin{alertblock}{Guiding question}
    If $p=0.18$ of tickets are late, how much does $\hat p$ move in a
    sample of $n=120$, and when is a bell the wrong picture?
  \end{alertblock}

  \vspace{0.6em}
  By the end of this lesson, you can:
  \begin{itemize}
    \item write $E[\hat p]=p$ and $\mathrm{SE}=\sqrt{p(1-p)/n}$;
    \item check the $np$ and $n(1-p)$ conditions;
    \item recover the SE from a seeded simulation;
    \item explain why $np<5$ breaks the normal model of $\hat p$.
  \end{itemize}

  \vfill
  \scriptsize All late-delivery flags used today are synthetic. No real company data are used.
\end{frame}

\begin{frame}{SE of $\hat p$ and the $np$ Check}
  \small
  Population $p=0.18$, sample size $n=120$:
  \[
    \mathrm{SE}(\hat p)=\sqrt{\frac{0.18\times 0.82}{120}}=0.035071.
  \]
  Counts:
  \[
    np=21.6,\qquad n(1-p)=98.4.
  \]
  Both exceed 5, so a normal picture of $\hat p$ is the working model.

  \vspace{0.3em}
  \begin{exampleblock}{Synthetic late-delivery flag}
    $\hat p$ is a sample statistic. $p$ is the unknown rate we study.
  \end{exampleblock}
\end{frame}

\begin{frame}[fragile]{Step 1: Formula SE}
  \begin{columns}[T]
    \begin{column}{0.42\textwidth}
      \small
      \begin{lstlisting}
se = sqrt(
  p * (1 - p) / n
)
np_count = n * p
      \end{lstlisting}

      \begin{block}{Verified output}
        $np=21.600000$\\
        $n(1-p)=98.400000$\\
        SE $=0.035071$
      \end{block}
    \end{column}
    \begin{column}{0.58\textwidth}
      \centering
      \includegraphics[width=\textwidth]{../figures/sampling_proportion_01_se_phat.png}
      \vspace{0.2em}
      \scriptsize Normal model of $\hat p$ from Step 1
    \end{column}
  \end{columns}
\end{frame}

\begin{frame}{Simulation Recovers the SE}
  \small
  5{,}000 seed-42 samples of $n=120$:
  \[
    \text{mean of }\hat p=0.179663,\qquad
    \widehat{\mathrm{SE}}=0.035079.
  \]
  Simulated $P(\hat p>0.22)=0.119800$ versus normal $0.127032$.
\end{frame}

\begin{frame}[fragile]{Step 2: 5{,}000 Values of $\hat p$}
  \begin{columns}[T]
    \begin{column}{0.40\textwidth}
      \small
      \begin{lstlisting}
successes = rng.binomial(
  n=120, p=0.18, size=5000
)
phat = successes / 120
      \end{lstlisting}

      \begin{block}{Verified output}
        Mean $=0.179663$\\
        Empirical SE $=0.035079$
      \end{block}
    \end{column}
    \begin{column}{0.60\textwidth}
      \centering
      \includegraphics[width=\textwidth]{../figures/sampling_proportion_02_simulation.png}
      \vspace{0.2em}
      \scriptsize Simulated $\hat p$ from Step 2
    \end{column}
  \end{columns}
\end{frame}

\begin{frame}{Step 3: $np<5$ Breaks the Bell}
  \begin{columns}[T]
    \begin{column}{0.42\textwidth}
      \small
      Now $n=20$, so $np=3.6<5$.

      \vspace{0.4em}
      Exact $P(\hat p=0)=0.018892$.\\
      Normal mass near 0 $=0.035594$.\\
      Skewness of $\hat p=0.430019$.

      \vspace{0.4em}
      \textbf{Limit:} $np<5$ makes the normal approximation fail.
      Lesson 26 studies point estimators.
    \end{column}
    \begin{column}{0.58\textwidth}
      \centering
      \includegraphics[width=\textwidth]{../figures/sampling_proportion_03_np_condition_limit.png}
      \vspace{0.2em}
      \scriptsize Discrete $\hat p$ versus illegal bell from Step 3
    \end{column}
  \end{columns}
\end{frame}

\begin{frame}{Concept Check}
  \begin{enumerate}
    \item Compute $\mathrm{SE}(\hat p)$ for $p=0.18$ and $n=120$.
    \item Why do $np=21.6$ and $n(1-p)=98.4$ matter?
    \item The empirical SE is $0.035079$. Does that replace $0.035071$?
    \item Why is $P(\hat p=0)=0.018892$ a problem for a normal model?
  \end{enumerate}

  \vspace{0.6em}
  \begin{block}{Connection to Lesson 26}
    The session can continue directly into unbiasedness, efficiency, and
    consistency of point estimators.
  \end{block}

  \vfill
  \scriptsize Source: OpenStax, \textit{Introductory Business Statistics 2e},
  Chapter 7, Section 7.3. CC BY-NC-SA.
\end{frame}

\appendix



\end{document}
